\begin{textsample}{POS Dim 6 – human – Score 4.00 – es/es\_ef\_anos\_iniciais\_ensino\_religioso\_1\_p1.txt}  \label{ex:f6_pos_005}
Ao longo da história da educação \textbf{brasileira}, o Ensino Religioso assumiu diferentes perspectivas teóricometodológicas, geralmente de viés confessional ou interconfessional.
A partir da década de 1980, as transformações socioculturais que provocaram mudanças paradigmáticas no campo educacional também impactaram no Ensino Religioso.
Em função dos promulgados ideais de democracia, inclusão social e educação \textbf{integral}, vários setores da sociedade civil passaram a reivindicar a abordagem do conhecimento religioso e o \textbf{reconhecimento} da diversidade religiosa no âmbito dos currículos escolares.

% matched lemmas: brasileiro, integral, reconhecimento
\end{textsample}
